\section{Related Work}

Sign language recognition has grown into an active research area in computer
vision. Early methods relied on handcrafted features and classical classifiers, but
convolutional neural networks (CNNs) largely replaced them because they learn
useful features directly from images~\cite{krizhevsky2012, he2016}. Since training a
large network from scratch needs a lot of data, most recent work uses transfer
learning, where a model pretrained on a large image dataset is adapted to the
gesture task. Lightweight backbones such as MobileNet~\cite{sandler2018mobilenetv2}
and ResNet~\cite{he2016} are common choices, because they run fast enough for
real-time use while still classifying accurately.

Most of this progress has gone to American Sign Language~\cite{rastgoo2021survey},
which now has many public datasets and recognition systems. Arabic Sign Language has
received far less attention~\cite{mohandes2014}, even though the need for Arabic
learning tools is clear. Several studies have applied CNNs to Arabic alphabet
recognition~\cite{latif2020arsl, mohandes2014} and report high accuracy. Two
limitations show up often: the systems are usually static classifiers with no
teaching component, and the accuracy is measured on a single controlled dataset,
which does not always reflect how the model behaves on new users or in different
conditions.

\subsection{Educational Sign Language Systems}
Some work looks beyond pure recognition at how computer vision can help people
learn. Tools that give immediate visual feedback tend to help beginners more than
passive video tutorials, since the learner can fix mistakes as they happen. Most of
these tools target ASL or European sign languages, so Arabic learners are left with
few options. Combining accurate recognition with a structured learning task is
still uncommon for Arabic, and that is the gap this work tries to fill.

\subsection{Transfer Learning for Limited-Data Domains}
Transfer learning is the usual choice when the target task has limited data. The
low-level features a network learns on natural images, such as edges and textures,
carry over well to hand images, so a pretrained backbone gives a strong starting
point. A common recipe is two-phase fine-tuning: first train only the new classifier
head while the backbone stays frozen, then unfreeze the whole network and continue
at a smaller learning rate. Freezing the backbone early keeps the untrained head
from corrupting the pretrained weights. We use this approach for the ArASL task.